\specialchapter{SELF-ADJOINTNESS CRITERIA}

We summarize here the criteria stated in the text that allow one to prove that certain partial operators are self-adjoint.

We fix a complex Hilbert space E and a partial operator u on E. This operator is self-adjoint in each of the following cases:

\begin{enumerate}
\def\labelenumi{(\roman{enumi})}
\item
  There exists \(v \in \mathcal{L}(\mathrm{E})\) injective and Hermitian such that \(u = v^{-1}\) . (Proposition 9 of IV, p. 239);
\item
  The operator \(u\) is symmetric and \(\mathrm{dom}(u) = \mathrm{E}\) (Proposition 10 of IV, p. 240);
\item
  The operator \(u\) is symmetric and induces a bijection of \(\mathrm{dom}(u)\) onto E (corollary of Proposition 10 of IV, p. 240);
\item
  There exists \(\lambda \in \mathbf{C}\) such that \(u + \lambda 1_{\mathrm{E}}\) and \(u + \overline{\lambda} 1_{\mathrm{E}}\) are surjective (Proposition 11 of IV, p. 240);
\item
  There exists a partial operator \(v\) , closed with dense domain, such that \(u = v^{*} \circ v\) (Proposition 12 of IV, p. 241);
\item
  The operator \(u\) is symmetric, closed, and its deficiency index is \((0,0)\) (corollary of Proposition 26 of IV, p. 257), that is, \(\mathrm{Ker}(u^{*} - i) = \mathrm{Ker}(u^{*} + i) = \{0\}\) ;
\item
  The space E is of the form \(\mathrm{L}^{2}(\mathrm{X},\mu)\) , where X is a locally compact topological space and \(\mu\) is a measure on X, and u is the operator of multiplication by a function \(\mu\) -measurable on X which is locally \(\mu\) -almost everywhere real-valued (corollary of Proposition 23);
\item
  There exists a normal partial operator v on E and a universally measurable function \(f \in \mathcal{L}_{\mathrm{u}}(\mathrm{Sp}(v))\) with real values such that \(u = f(v)\) (Corollary 1 of IV, p. 273);
\item
  There exists a partial Hermitian form q on E such that u is the partial operator representing q (Proposition 14 of IV, p. 293);
\item
  One has \(u = v + w\) where v is a self-adjoint partial operator and w is a symmetric partial operator bounded relative to v satisfying \(\|w\|_{v} < 1\) (Theorem 5 of IV, p. 306);
\item
  The operator \(u\) is the infinitesimal generator of a unitary representation \(\varrho\) of \(\mathbf{R}\) in E (Theorem 1 of V, p. 428).
\end{enumerate}
